% Work/MAT | 1. rocnik | 36: Vedecky a technicky zapis cisel
% Puvodni rozpis mel vedecky zapis ve 29. hodine, novy plan ho presunul na 36. hodinu.
% Terminologie: vedecky (normalizovany) tvar 1<=|a|<10, n cele;
% technicky / inzenyrsky tvar 1<=|a|<1000, exponent nasobkem tri.
% Zdroje: ucebnice Mocniny a Odmocniny.tex (Mocniny desitky a vedecky
% zapis cisla); Realisticky.cz 010608 Exponencialni tvar cisla,
% 010609 Vypocty s mocninami. Texty a priklady jsou originalni.
\documentclass[11pt,a4paper]{article}
% Spolecny styl pro ucitelske pripravy MAT (Work/MAT)
% Verze 1.1, 2026-10-23
\NeedsTeXFormat{LaTeX2e}
\ProvidesPackage{priprava}[2026/10/23 v1.1 Sablona ucitelskych priprav MAT]

% Zakladni typografie a cestina
\RequirePackage[T1]{fontenc}
\RequirePackage[utf8]{inputenc}
\RequirePackage[czech]{babel}
\RequirePackage{lmodern}
\RequirePackage{microtype}
\RequirePackage{amsmath,amssymb}
\RequirePackage{xcolor}
\RequirePackage{enumitem}
\RequirePackage[most]{tcolorbox}
\RequirePackage{tabularx,booktabs,array}
\RequirePackage{fancyhdr}
\RequirePackage{needspace}
\RequirePackage{geometry}
\RequirePackage[hidelinks]{hyperref}

\geometry{a4paper,top=17mm,bottom=18mm,left=23mm,right=23mm}
\definecolor{notegray}{gray}{0.39}
\setlength{\parindent}{0pt}
\setlength{\parskip}{2.5pt plus 1pt minus 0.5pt}
\setlist[itemize]{topsep=2pt,itemsep=1pt,parsep=0pt,leftmargin=1.45em}
\setlist[enumerate]{topsep=2pt,itemsep=2pt,parsep=0pt,leftmargin=1.65em}
\setcounter{secnumdepth}{0}
\renewcommand{\arraystretch}{1.12}

% Zahlavi: vlevo rocnik | globalni tema, vpravo cislo: nazev hodiny
% Zapati zustava zcela prazdne.
\pagestyle{fancy}
\fancyhf{}
\renewcommand{\headrulewidth}{0.35pt}
\renewcommand{\footrulewidth}{0pt}
\setlength{\headheight}{14pt}
\newcommand{\nastavPripravu}[4]{%
  \fancyhead[L]{\small #1\enspace\textbar\enspace #2}%
  \fancyhead[R]{\small #3: #4}%
}

% Minimalisticke cernobile boxy. Kazda definice zvlast.
\tcbset{commonlesson/.style={enhanced,breakable=false,colback=white,colframe=black,
  colbacktitle=white,coltitle=black,boxrule=0.45pt,
  arc=0.8mm,outer arc=0.8mm,left=3mm,right=3mm,top=1.6mm,bottom=1.6mm,
  toptitle=1mm,bottomtitle=0.7mm,fonttitle=\bfseries,
  attach boxed title to top left={xshift=2.5mm,yshift=-1.15mm},
  boxed title style={colback=white,colframe=white,boxrule=0pt,left=1mm,right=1mm}}}
\newtcolorbox{zadaniBox}[1][]{commonlesson,title={Zadání},#1}
\newtcolorbox{prikladBox}[1][]{commonlesson,title={Příklad},#1}
\newtcolorbox{definiceBox}[1][]{commonlesson,title={Definice},#1}

% Spolecny box pro zadani a jeho reseni, bez mezery mezi castmi.
% Pouziti:
%   \begin{prikladReseniBox}
%     Zadani ulohy.
%     \tcblower
%     \begin{align*} ... \end{align*}
%   \end{prikladReseniBox}
% Spodni cast ma jemny sedy podklad a delici linku; lze v ni sazet
% vice radku matematiky. Obe casti maji jen jeden spolecny ramecek.
\tcbset{solutionlesson/.style={
  commonlesson,
  enhanced jigsaw,
  breakable,
  colbacklower=black!3,
  segmentation style={solid,draw=black!25,line width=0.35pt},
  fontlower=\small\color{notegray},
  before lower={\noindent\textbf{Řešení:}\enspace},
  leftlower=3mm,rightlower=3mm
}}
\newtcolorbox{prikladReseniBox}[1][]{solutionlesson,title={Příklad},#1}
\newtcolorbox{zadaniReseniBox}[1][]{solutionlesson,title={Zadání},#1}

% Pomucky pro autora pripravy: sedive poznamky; \reseni pro starsi pripravy
\newcommand{\poznamka}[1]{\par\begingroup\small\color{notegray}\itshape
  \textnormal{\textbf{Poznámka pro učitele.}}\enspace #1\par\endgroup}
\newcommand{\reseni}[1]{\par\begingroup\small\color{notegray}
  \textbf{Očekávané řešení:}\enspace #1\par\endgroup}
\newcommand{\blok}[1]{\needspace{5\baselineskip}\subsection*{#1}}
\newcommand{\vyrok}[1]{\uv{#1}}

\usepackage{multicol}
\nastavPripravu{1. ročník}{Mocniny a odmocniny}{36}{Vědecký a technický zápis čísel}
\begin{document}
\poznamka{\textbf{45minutové jádro:} 4 min motivace a řády;
8 min definice a Příklad 1; 8 min Samostatná práce 1 a)--f);
7 min Příklad 2 a Samostatná práce 2 a)--d);
8 min Příklad 3 (operace) a Samostatná práce 3 a)--c);
7 min Příklad 4, vybrané předpony jednotek; 3 min shrnutí a domácí práce.
Technický zápis s násobky tří lze projít stručně; rozšiřující příklady,
informatické jednotky jsou nadstavba.
Nenech se svést k předčasnému zavádění exponenciální funkce.}

\blok{Proč psát mocniny desítky?}
Porovnej zápisy $4\,700\,000\,000$ a $4{,}7\cdot10^9$.
Také $0{,}0000047=4{,}7\cdot10^{-6}$. Mocnina desítky
ukazuje \textbf{řád čísla} a usnadňuje srovnávání velmi velkých
i velmi malých hodnot. Záporný exponent nepředstavuje záporné číslo:
$10^{-3}=0{,}001>0$.

\begin{center}\small
\begin{tabular}{@{}c|rrrrrrr@{}}
\toprule
Exponent $n$ & $3$ & $2$ & $1$ & $0$ & $-1$ & $-2$ & $-3$\\\midrule
$10^n$ & $1000$ & $100$ & $10$ & $1$ & $0{,}1$ & $0{,}01$ & $0{,}001$\\
\bottomrule
\end{tabular}
\end{center}
\poznamka{Nech žáky bez kalkulačky říct, jak se mění hodnota,
když exponent zvýšíme či snížíme o jedna. Každý krok doprava v tabulce
dělí hodnotu deseti.}

\blok{1. Vědecký zápis čísla}
\begin{definiceBox}[title={Definice: vědecký (normalizovaný exponenciální) zápis}]
Každé \textbf{nenulové reálné číslo} $x$ lze jednoznačně zapsat ve tvaru
\[
\boxed{x=a\cdot10^n,\qquad 1\le |a|<10,\quad n\in\mathbb Z.}
\]
Číslo $a$ je \textbf{koeficient} (znaménko zůstává u něho),
$n$ je \textbf{řádový exponent}. Číslo $0$ zapisujeme jednoduše jako $0$.
\end{definiceBox}
\textbf{Postup:} desetinnou čárku posuň tak, aby před ní byla právě
jedna nenulová číslice. Posuneš-li ji \emph{doleva} o $k$ míst,
exponent je $+k$; posuneš-li ji \emph{doprava}, exponent je $-k$.

\begin{prikladReseniBox}[breakable=false,title={Příklad 1: z desetinného zápisu do vědeckého a zpět}]
Převeď na vědecký zápis $5\,280\,000$,
$0{,}000072$ a $-430\,000$.
Na běžný desetinný zápis převeď $6{,}2\cdot10^{-4}$.
\tcblower
\[
5\,280\,000=\boxed{5{,}28\cdot10^6},\qquad
0{,}000072=\boxed{7{,}2\cdot10^{-5}},
\]
\[
-430\,000=\boxed{-4{,}3\cdot10^5},\qquad
6{,}2\cdot10^{-4}=\boxed{0{,}00062}.
\]
Všimni si, že znaménko čísla a znaménko exponentu vyjadřují různé věci.
\end{prikladReseniBox}

\begin{zadaniReseniBox}[breakable=false,title={Samostatná práce 1: tam a zpět}]
\begin{multicols}{2}
\begin{enumerate}[label=\alph*)]
\item $7\,300\,000$;
\item $0{,}0000049$;
\item $-0{,}0082$;
\item $91\,500$;
\item $3{,}7\cdot10^{-3}$;
\item $-8{,}04\cdot10^5$;
\item $0{,}000000031$;
\item $4{,}2\cdot10^7$.
\end{enumerate}
\end{multicols}
U a)--d), g) zapiš vědecký tvar; u e), f), h) běžné číslo.
\tcblower
\begin{multicols}{2}
\begin{enumerate}[label=\alph*)]
\item $7{,}3\cdot10^6$;
\item $4{,}9\cdot10^{-6}$;
\item $-8{,}2\cdot10^{-3}$;
\item $9{,}15\cdot10^4$;
\item $0{,}0037$;
\item $-804\,000$;
\item $3{,}1\cdot10^{-8}$;
\item $42\,000\,000$.
\end{enumerate}
\end{multicols}
\end{zadaniReseniBox}

\blok{2. Zápis na kalkulačce a kontrola řádu}
Kalkulačka může ukázat \texttt{2.5E-4} nebo \texttt{2.5e-4}.
To znamená $2{,}5\cdot10^{-4}$, \textbf{nikoli} násobení Eulerovým číslem $e$.
Tlačítko \texttt{EXP} či \texttt{EE} vkládá \emph{řád desítky}; jeho přesný název
závisí na kalkulačce. Například \texttt{6.2 EXP 7} značí $6{,}2\cdot10^7$.

\begin{prikladReseniBox}[breakable=false,title={Příklad 2: co je a není správně normalizované?}]
Posuď zápisy $45\cdot10^4$, $4{,}5\cdot10^5$,
$0{,}45\cdot10^6$ a $-4{,}5\cdot10^5$.
\tcblower
První tři zápisy vyjadřují stejnou \emph{kladnou} hodnotu
$450\,000$, ale v normalizovaném vědeckém tvaru je pouze
$\boxed{4{,}5\cdot10^5}$: právě zde $1\le|a|<10$.
Poslední zápis je správný vědecký tvar čísla $-450\,000$.
\end{prikladReseniBox}

\begin{zadaniReseniBox}[breakable=false,title={Samostatná práce 2: oprav nenormalizovaný zápis}]
Přepiš do vědeckého tvaru:
\begin{multicols}{2}
\begin{enumerate}[label=\alph*)]
\item $36\cdot10^5$;
\item $0{,}005\cdot10^3$;
\item $-250\cdot10^{-4}$;
\item $0{,}72\cdot10^{-7}$;
\item $100\cdot10^0$;
\item $-0{,}003\cdot10^{-2}$.
\end{enumerate}
\end{multicols}
\tcblower
\begin{multicols}{2}
\begin{enumerate}[label=\alph*)]
\item $3{,}6\cdot10^6$;
\item $5\cdot10^0=5$;
\item $-2{,}5\cdot10^{-2}$;
\item $7{,}2\cdot10^{-8}$;
\item $1\cdot10^2$;
\item $-3\cdot10^{-5}$.
\end{enumerate}
\end{multicols}
\end{zadaniReseniBox}

\blok{3. Počítání s čísly ve vědeckém tvaru}
Při násobení a dělení použij známá pravidla pro mocniny. Při
\textbf{sčítání a odčítání} však musíš nejdřív sjednotit exponenty;
koeficienty nemůžeš prostě sečíst a exponenty také.
\begin{prikladReseniBox}[breakable=false,title={Příklad 3: násobení a dělení bez velkých čísel}]
Vypočti a napiš výsledek ve vědeckém tvaru:
\[
(3{,}2\cdot10^{-4})(1{,}5\cdot10^6),\qquad
\frac{6\cdot10^8}{2\cdot10^3}.
\]
\tcblower
\[
(3{,}2\cdot1{,}5)10^{-4+6}=4{,}8\cdot10^2,\qquad
\frac62\cdot10^{8-3}=3\cdot10^5.
\]
\textbf{Kontrola:} $4{,}8\cdot10^2=480$.
\end{prikladReseniBox}

\begin{zadaniReseniBox}[breakable=false,title={Samostatná práce 3: výpočty a řád výsledku}]
Vypočti ve vědeckém tvaru:
\begin{enumerate}[label=\alph*)]
\item $(2\cdot10^5)(4\cdot10^{-3})$;
\item $\displaystyle\frac{8{,}4\cdot10^{-3}}{2{,}1\cdot10^{-5}}$;
\item $(7{,}2\cdot10^4)(5\cdot10^{-3})$;
\item $6\cdot10^4+3\cdot10^3$;
\item $9\cdot10^{-2}-4\cdot10^{-3}$;
\item $\displaystyle\frac{3\cdot10^6}{6\cdot10^{-2}}$.
\end{enumerate}
\tcblower
\begin{enumerate}[label=\alph*)]
\item $8\cdot10^2$;
\item $(8{,}4/2{,}1)10^{-3+5}=4\cdot10^2$;
\item $36\cdot10^1=3{,}6\cdot10^2$;
\item $(6+0{,}3)10^4=6{,}3\cdot10^4$;
\item $(9-0{,}4)10^{-2}=8{,}6\cdot10^{-2}$;
\item $(3/6)10^8=5\cdot10^7$.
\end{enumerate}
\end{zadaniReseniBox}
\poznamka{Do jádra stačí a)--c). U d), e) ověř, zda studenti chápou,
proč nelze u sčítání použít exponent $m+n$.}

\blok{4. Technický (inženýrský) zápis a předpony jednotek}
\begin{definiceBox}[title={Definice: technický zápis}]
V \textbf{technickém (inženýrském) zápisu} má nenulové číslo tvar
\[
\boxed{x=b\cdot10^{3k},\qquad1\le|b|<1000,\quad k\in\mathbb Z.}
\]
Exponent je \textbf{násobkem tří}. To se hodí k předponám
jednotek, které postupují po tisícinách a tisícinásobcích.
\end{definiceBox}
\begin{center}\small
\begin{tabularx}{\linewidth}{@{}>{\raggedright\arraybackslash}p{19mm}>{\raggedright\arraybackslash}p{18mm}>{\raggedright\arraybackslash}p{23mm}X@{}}
\toprule
Mocnina & Předpona & Značka & Příklad\\\midrule
$10^9$ & giga & G & $3{,}2\,\mathrm{GHz}=3{,}2\cdot10^9\,\mathrm{Hz}$\\
$10^6$ & mega & M & $2\,\mathrm{MW}=2\cdot10^6\,\mathrm W$\\
$10^3$ & kilo & k & $5\,\mathrm{km}=5\cdot10^3\,\mathrm m$\\
$10^{-3}$ & mili & m & $4\,\mathrm{mm}=4\cdot10^{-3}\,\mathrm m$\\
$10^{-6}$ & mikro & $\mu$ & $7\,\mathrm{\mu s}=7\cdot10^{-6}\,\mathrm s$\\
$10^{-9}$ & nano & n & $9\,\mathrm{nm}=9\cdot10^{-9}\,\mathrm m$\\
\bottomrule
\end{tabularx}
\end{center}

\begin{prikladReseniBox}[breakable=false,title={Příklad 4: jedna hodnota, dva užitečné zápisy}]
Zapiš $0{,}00047\,\mathrm m$ vědecky, technicky a pomocí vhodné předpony.
\tcblower
Vědecky $4{,}7\cdot10^{-4}\,\mathrm m$,
technicky $470\cdot10^{-6}\,\mathrm m$,
tedy $\boxed{470\,\mathrm{\mu m}}$.
Obě podoby vyjadřují tutéž délku; měníme pouze způsob zápisu.
\end{prikladReseniBox}

\begin{zadaniReseniBox}[breakable=false,title={Samostatná práce 4: technické veličiny}]
\begin{enumerate}[label=\alph*)]
\item $3\,700\,000\,\mathrm{Hz}$ vyjádři v MHz;
\item $0{,}00000082\,\mathrm m$ vyjádři v nm;
\item $0{,}0000000042\,\mathrm s$ vyjádři v ns;
\item $4\,500\,000\,\mathrm W$ vyjádři v MW.
\end{enumerate}
\tcblower
\begin{enumerate}[label=\alph*)]
\item $3{,}7\,\mathrm{MHz}$;
\item $820\,\mathrm{nm}$;
\item $4{,}2\,\mathrm{ns}$;
\item $4{,}5\,\mathrm{MW}$.
\end{enumerate}
\end{zadaniReseniBox}

\blok{Rozšíření pro technicky zaměřené žáky}
\begin{zadaniReseniBox}[breakable=false,title={Rozšiřující příklad 1: kontrola jednotek a řádu}]
\begin{enumerate}[label=\alph*)]
\item Čtvercová destička má stranu $4\cdot10^{-3}\,\mathrm m$. Urči obsah v $\mathrm{m^2}$.
\item Kostka má hranu $2\cdot10^{-3}\,\mathrm m$. Urči objem v $\mathrm{m^3}$.
\item Rozliš: proč je $1\,\mathrm{kB}=1000\,\mathrm B$, ale $1\,\mathrm{KiB}=1024\,\mathrm B$?
\end{enumerate}
\tcblower
\begin{enumerate}[label=\alph*)]
\item $(4\cdot10^{-3})^2=16\cdot10^{-6}=\boxed{1{,}6\cdot10^{-5}\,\mathrm{m^2}}$;
\item $(2\cdot10^{-3})^3=8\cdot10^{-9}\,\mathrm{m^3}$;
\item Předpona k (kilo) vychází z desítkové soustavy, Ki (kibi) ze dvojkové.
\end{enumerate}
\end{zadaniReseniBox}
\begin{zadaniReseniBox}[breakable=false,title={Rozšiřující příklad 2: časté omyly}]
Rozhodni, zda tvrzení platí; chybu oprav:
\begin{enumerate}[label=\alph*)]
\item $5{,}2\cdot10^{-7}$ je záporné číslo;
\item $3{,}4\cdot10^7+2\cdot10^6=5{,}4\cdot10^{13}$;
\item $12\cdot10^3$ je normalizovaný vědecký zápis.
\end{enumerate}
\tcblower
\begin{enumerate}[label=\alph*)]
\item Ne; hodnota je kladná;
\item Ne; $3{,}4\cdot10^7+0{,}2\cdot10^7=3{,}6\cdot10^7$;
\item Ne; vědecky $1{,}2\cdot10^4$ (technicky je $12\cdot10^3$ správně).
\end{enumerate}
\end{zadaniReseniBox}

\blok{Domácí procvičování}
\begin{zadaniReseniBox}[breakable=false,title={Domácí úkol 1: vědecký a technický zápis}]
\begin{enumerate}[label=\alph*)]
\item Zapiš vědecky $89\,000\,000$;
\item zapiš vědecky $0{,}00000056$;
\item napiš desetinně $-7{,}3\cdot10^{-4}$;
\item uprav $(2{,}5\cdot10^6)(4\cdot10^{-8})$;
\item uprav $\dfrac{9\cdot10^{-5}}{3\cdot10^2}$;
\item vypočítej $4\cdot10^5+7\cdot10^4$;
\item přepiš $6{,}4\cdot10^{-5}$ do technického tvaru;
\item vyjádři $0{,}000009\,\mathrm m$ v $\mathrm{\mu m}$.
\end{enumerate}
\tcblower
\begin{enumerate}[label=\alph*)]
\item $8{,}9\cdot10^7$;
\item $5{,}6\cdot10^{-7}$;
\item $-0{,}00073$;
\item $10\cdot10^{-2}=1\cdot10^{-1}$;
\item $3\cdot10^{-7}$;
\item $4{,}7\cdot10^5$;
\item $64\cdot10^{-6}$;
\item $9\,\mathrm{\mu m}$.
\end{enumerate}
\end{zadaniReseniBox}

\poznamka{\textbf{Výstupní otázka:} Jaký je rozdíl mezi vědeckým tvarem
$1{,}8\cdot10^4$ a technickým tvarem $18\cdot10^3$?
Číslo je stejné, ale technický exponent je násobkem tří.
\textbf{Organizace rozpisu:} hodiny 35 a 36 jsou teď pracovními hodinami;
původní termín čtvrtletní písemné práce na 32. hodině byl posunut a je
nutné mu určit nové místo v celkovém časovém plánu.}
\end{document}
